\SourcePage{789}{792}
\renewcommand{\thesection}{\alph{section}}
\setcounter{section}{4}
\section{The hypergeometric function}

Within the circle $|z|<1$, the \emph{hypergeometric function} is defined by
the series
\begin{equation}
 F(\alpha,\beta,\gamma,z)=1+\frac{\alpha\beta}{\gamma}\frac z{1!}
 +\frac{\alpha(\alpha+1)\beta(\beta+1)}{\gamma(\gamma+1)}
 \frac{z^2}{2!}+\ldots,
 \tag{e.1}
\end{equation}
and for $|z|>1$ it is obtained by analytic continuation of this series (see
(e.6)). The hypergeometric function is one particular integral of the
differential equation
\begin{equation}
 z(1-z)u''+[\gamma-(\alpha+\beta+1)z]u'-\alpha\beta u=0.
 \tag{e.2}
\end{equation}
The parameters $\alpha$ and $\beta$ are arbitrary, while
$\gamma\ne0,-1,-2,\ldots$. Evidently, $F(\alpha,\beta,\gamma,z)$ is symmetric
in the parameters $\alpha$ and $\beta$\footnote{The confluent hypergeometric
function is obtained from $F(\alpha,\beta,\gamma,z)$ by the limiting
transition
$F(\alpha,\gamma,z)=\lim_{\beta\to\infty}
F(\alpha,\beta,\gamma,z/\beta)$. The notations
${}_2F_1(\alpha,\beta,\gamma,z)$ and ${}_1F_1(\alpha,\gamma,z)$ are also used
in the literature for the hypergeometric and confluent hypergeometric
functions, respectively. The indices to the left and right of $F$ give the
numbers of parameters appearing in the numerators and denominators,
respectively, of the terms of the series.}.

A second independent solution of (e.2) is
\[
 z^{1-\gamma}F(\beta-\gamma+1,\alpha-\gamma+1,2-\gamma,z);
\]
it has a singular point at $z=0$.

For reference, we give here a number of relations satisfied by the
hypergeometric function.

For all $z$, provided $\operatorname{Re}(\gamma-\alpha)>0$, the function
$F(\alpha,\beta,\gamma,z)$ can be represented by the integral
\begin{equation}
 \begin{split}
 F(\alpha,\beta,\gamma,z)
 &=-\frac1{2\pi i}
 \frac{\Gamma(1-\alpha)\Gamma(\gamma)}{\Gamma(\gamma-\alpha)}\\
 &\hspace{1.5em}\times\oint_{C'}(-t)^{\alpha-1}
 (1-t)^{\gamma-\alpha-1}(1-tz)^{-\beta}\,dt,
 \end{split}
 \tag{e.3}
\end{equation}
taken over the contour $C'$ shown in Fig.~57. Direct substitution readily
verifies that this integral does indeed satisfy (e.2); the constant factor has
been chosen so that its value at $z=0$ is unity.

The substitution
\[
 u=(1-z)^{\gamma-\alpha-\beta}u_1
\]

\SourcePage{790}{793}
in (e.2) gives an equation of the same form, with the parameters
$\gamma-\alpha$, $\gamma-\beta$, $\gamma$ in place of $\alpha$, $\beta$,
$\gamma$, respectively. It follows that
\begin{equation}
 F(\alpha,\beta,\gamma,z)=(1-z)^{\gamma-\alpha-\beta}
 F(\gamma-\alpha,\gamma-\beta,\gamma,z)
 \tag{e.4}
\end{equation}
(both sides obey the same equation and have equal values at $z=0$).

Making the substitution $t\to t/(1-z+zt)$ in the integral (e.3) gives the
following relation between hypergeometric functions of the variables $z$ and
$z/(z-1)$:
\begin{equation}
 F(\alpha,\beta,\gamma,z)=(1-z)^{-\alpha}
 F\left(\alpha,\gamma-\beta,\gamma,\frac z{z-1}\right).
 \tag{e.5}
\end{equation}
The value of the multivalued expression $(1-z)^{-\alpha}$ in this formula (and
of analogous expressions in all the formulas below) is fixed by requiring the
complex quantity being raised to a power to be taken with the argument of
smallest absolute value.

We next state, without proof, an important formula relating hypergeometric
functions of the variables $z$ and $1/z$:
\begin{equation}
 \begin{split}
 F(\alpha,\beta,\gamma,z)
 &=\frac{\Gamma(\gamma)\Gamma(\beta-\alpha)}
 {\Gamma(\beta)\Gamma(\gamma-\alpha)}(-z)^{-\alpha}
 F\left(\alpha,\alpha+1-\gamma,\alpha+1-\beta,\frac1z\right)\\
 &\hspace{1em}+\frac{\Gamma(\gamma)\Gamma(\alpha-\beta)}
 {\Gamma(\alpha)\Gamma(\gamma-\beta)}(-z)^{-\beta}
 F\left(\beta,\beta+1-\gamma,\beta+1-\alpha,\frac1z\right).
 \end{split}
 \tag{e.6}
\end{equation}
This formula expresses $F(\alpha,\beta,\gamma,z)$ as a series convergent for
$|z|>1$; it therefore supplies the analytic continuation of the original
series (e.1).

The formula
\begin{equation}
 \begin{split}
 F(\alpha,\beta,\gamma,z)
 &=\frac{\Gamma(\gamma)\Gamma(\gamma-\alpha-\beta)}
 {\Gamma(\gamma-\alpha)\Gamma(\gamma-\beta)}
 F(\alpha,\beta,\alpha+\beta+1-\gamma,1-z)\\
 &\hspace{1em}+\frac{\Gamma(\gamma)\Gamma(\alpha+\beta-\gamma)}
 {\Gamma(\alpha)\Gamma(\beta)}(1-z)^{\gamma-\alpha-\beta}\\
 &\hspace{3em}\times
 F(\gamma-\alpha,\gamma-\beta,\gamma+1-\alpha-\beta,1-z)
 \end{split}
 \tag{e.7}
\end{equation}
relates hypergeometric functions of $z$ and $1-z$ (we give this relation also
without proof). Combining (e.7) with (e.6), we obtain

\SourcePage{791}{794}
the relations
\begin{equation}
 \begin{split}
 F(\alpha,\beta,\gamma,z)
 &=\frac{\Gamma(\gamma)\Gamma(\beta-\alpha)}
 {\Gamma(\beta)\Gamma(\gamma-\alpha)}(1-z)^{-\alpha}
 F\left(\alpha,\gamma-\beta,\alpha+1-\beta,\frac1{1-z}\right)\\
 &\hspace{1em}+\frac{\Gamma(\gamma)\Gamma(\alpha-\beta)}
 {\Gamma(\alpha)\Gamma(\gamma-\beta)}(1-z)^{-\beta}
 F\left(\beta,\gamma-\alpha,\beta+1-\alpha,\frac1{1-z}\right),
 \end{split}
 \tag{e.8}
\end{equation}
\begin{equation}
 \begin{split}
 F(\alpha,\beta,\gamma,z)
 &=\frac{\Gamma(\gamma)\Gamma(\gamma-\alpha-\beta)}
 {\Gamma(\gamma-\beta)\Gamma(\gamma-\alpha)}z^{-\alpha}
 F\left(\alpha,\alpha+1-\gamma,\alpha+\beta+1-\gamma,
 \frac{z-1}{z}\right)\\
 &\hspace{1em}+\frac{\Gamma(\gamma)\Gamma(\alpha+\beta-\gamma)}
 {\Gamma(\alpha)\Gamma(\beta)}(1-z)^{\gamma-\alpha-\beta}z^{\beta-\gamma}\\
 &\hspace{3em}\times
 F\left(1-\beta,\gamma-\beta,\gamma+1-\alpha-\beta,
 \frac{z-1}{z}\right).
 \end{split}
 \tag{e.9}
\end{equation}
Each term in the sums on the right-hand sides of (e.6)--(e.9) is itself a
solution of the hypergeometric equation.

If $\alpha$ (or $\beta$) is a negative integer (or zero), $\alpha=-n$, the
hypergeometric function becomes a polynomial of degree $n$ and may be written
as
\begin{equation}
 F(-n,\beta,\gamma,z)=
 \frac{z^{1-\gamma}(1-z)^{\gamma+n-\beta}}
 {\gamma(\gamma+1)\cdots(\gamma+n-1)}
 \frac{d^n}{dz^n}\left[z^{\gamma+n-1}(1-z)^{\beta-\gamma}\right].
 \tag{e.10}
\end{equation}
Apart from a constant factor, these polynomials are the \emph{Jacobi
polynomials}, defined by
\begin{equation}
 \begin{split}
 P_n^{(a,b)}(z)
 &=\frac{(a+1)(a+2)\cdots(a+n)}{n!}
 F\left(-n,a+b+n+1,a+1,\frac{1-z}{2}\right)\\
 &=\frac{(-1)^n}{2^nn!}(1-z)^{-a}(1+z)^{-b}
 \frac{d^n}{dz^n}\left[(1-z)^{a+n}(1+z)^{b+n}\right].
 \end{split}
 \tag{e.11}
\end{equation}
For $a=b=0$, the Jacobi polynomials coincide with the Legendre polynomials.
For $n=0$, $P_0^{(a,b)}=1$.
